\documentclass[10pt,a4paper]{amsart}
\usepackage{amsmath,amssymb,amsthm,graphicx}
\usepackage{verbatim}
\newenvironment{reference}{\comment}{\endcomment}
%\newenvironment{reference}{}{}
\newenvironment{slogan}{\comment}{\endcomment}
\newenvironment{history}{\comment}{\endcomment}

% For commutative diagrams we use Xy-pic
\usepackage[all]{xy}

% We use 2cell for 2-commutative diagrams.
\xyoption{2cell}
% Theorem environments.
%
\theoremstyle{plain}
\newtheorem{theorem}[subsection]{Theorem}
\newtheorem{proposition}[subsection]{Proposition}
\newtheorem{lemma}[subsection]{Lemma}

\theoremstyle{definition}
\newtheorem{definition}[subsection]{Definition}
\newtheorem{example}[subsection]{Example}
\newtheorem{exercise}[subsection]{Exercise}
\newtheorem{situation}[subsection]{Situation}

\theoremstyle{remark}
\newtheorem{remark}[subsection]{Remark}
\newtheorem{remarks}[subsection]{Remarks}

\numberwithin{equation}{subsection}

% Macros
%
\def\lim{\mathop{\mathrm{lim}}\nolimits}
\def\colim{\mathop{\mathrm{colim}}\nolimits}
\def\Spec{\mathop{\mathrm{Spec}}}
\def\Hom{\mathop{\mathrm{Hom}}\nolimits}
\def\Ext{\mathop{\mathrm{Ext}}\nolimits}
\def\SheafHom{\mathop{\mathcal{H}\!\mathit{om}}\nolimits}
\def\SheafExt{\mathop{\mathcal{E}\!\mathit{xt}}\nolimits}
\def\Sch{\mathit{Sch}}
\def\Mor{\mathop{\mathrm{Mor}}\nolimits}
\def\Ob{\mathop{\mathrm{Ob}}\nolimits}
\def\Sh{\mathop{\mathit{Sh}}\nolimits}
\def\NL{\mathop{N\!L}\nolimits}
\def\CH{\mathop{\mathrm{CH}}\nolimits}
\def\proetale{{pro\text{-}\acute{e}tale}}
\def\etale{{\acute{e}tale}}
\def\QCoh{\mathit{QCoh}}
\def\Ker{\mathop{\mathrm{Ker}}}
\def\Im{\mathop{\mathrm{Im}}}
\def\Coker{\mathop{\mathrm{Coker}}}
\def\Coim{\mathop{\mathrm{Coim}}}

% Boxtimes
%
\DeclareMathSymbol{\boxtimes}{\mathbin}{AMSa}{"02}

%
% Macros for moduli stacks/spaces
%
\def\QCohstack{\mathcal{QC}\!\mathit{oh}}
\def\Cohstack{\mathcal{C}\!\mathit{oh}}
\def\Spacesstack{\mathcal{S}\!\mathit{paces}}
\def\Quotfunctor{\mathrm{Quot}}
\def\Hilbfunctor{\mathrm{Hilb}}
\def\Curvesstack{\mathcal{C}\!\mathit{urves}}
\def\Polarizedstack{\mathcal{P}\!\mathit{olarized}}
\def\Complexesstack{\mathcal{C}\!\mathit{omplexes}}
% \Pic is the operator that assigns to X its picard group, usage \Pic(X)
% \Picardstack_{X/B} denotes the Picard stack of X over B
% \Picardfunctor_{X/B} denotes the Picard functor of X over B
\def\Pic{\mathop{\mathrm{Pic}}\nolimits}
\def\Picardstack{\mathcal{P}\!\mathit{ic}}
\def\Picardfunctor{\mathrm{Pic}}
\def\Deformationcategory{\mathcal{D}\!\mathit{ef}}
\title{Subextensions of a radical extension of degree e are generated by powers of the radical}
\author{The Stacks Project Authors}
\date{stacks-project:fields.tex:local-snapshot-20261012; GNU FDL 1.2 or later}
\begin{document}
\maketitle
\noindent\textit{Editorial note.} All delivered mathematical text is the author's own, copied from the pinned Stacks Project chapter file fields.tex; no statement or proof step was rewritten, reordered or removed.\par\medskip
\noindent\textit{Editorial note.} The counted proof is the author's complete argument: The constant-term argument survives repeated roots: each root beta of P satisfies beta\textasciicircum{}e = a, so beta/zeta-free form beta = zeta alpha with zeta\textasciicircum{}e = 1, and zeta = prod zeta\_i satisfies zeta\textasciicircum{}e = 1.\par\medskip
\noindent\textit{Editorial note.} The label lemma-subfields-kummer names the counted statement, and the proof environment immediately following it in the source is delivered verbatim as its proof.\par\medskip
\noindent\textit{Editorial note.} The counted statement is carried by the author's own lemma environment in the chapter source.\par\medskip
\noindent\textit{Allowed standard prerequisites (editorial).}\begin{itemize}
\item Statement hypothesis: L = K(alpha) of degree e with a = alpha\textasciicircum{}e in K, and every e-th root of unity in L lies in K
\item Same chapter: multiplicativity of degrees in towers (lemma-multiplicativity-degrees)
\end{itemize}\medskip
\begin{lemma}
\label{lemma-subfields-kummer}
\begin{reference}
\cite[Theorem 5.2]{Radical}
\end{reference}
Let $K$ be a field. Let $L/K$ be a finite extension of degree $e$
which is generated by an element $\alpha$ with $a = \alpha^e \in K$.
If every $e$th root of unity in $L$ is contained in $K$, then 
any sub extension $L/L'/K$ is generated by $\alpha^d$ for some $d | e$.
\end{lemma}

\begin{proof}
Observe that for $d | e$ the subfield $K(\alpha^d)$ has
$[K(\alpha^d) : K] = e/d$ and $[L : K(\alpha^d)] = d$.
Let $L/L'/K$ be a subextension. Say $d = [L : L']$. If $\alpha^d \in L'$,
then we have $L' = K(\alpha^d)$ for degree reasons.
Let $P \in L'[x]$ be the minimal polynomial of $\alpha$ over $L'$.
Then $P$ divides $x^e - a$ and $P$ has degree $d$.
Let us write
$$
x^e - a = \prod\nolimits_{i = 1, \ldots, e} (x - \zeta_i \alpha)
$$
in a splitting field of $x^e - a$ over $L$. The $\zeta_i$ are $e$th roots
of unity and after renumbering we have
$$
P = \prod\nolimits_{i = 1, \ldots, d} (x - \zeta_i \alpha)
$$
The constant term of $P$ is equal to
$$
c = (\prod\nolimits_{i = 1, \ldots, d} \zeta_i) \alpha^d
$$
and is in $L' \subset L$. Since $\alpha \in L$ this implies that
$\zeta = \prod_{i = 1, \ldots, d} \zeta_i$ is in $L$ and hence in $K$ by
our assumption. Thus $\alpha^d = \zeta^{-1}c \in L'$ and we conclude.
\end{proof}

\end{document}
